\chapter{Problem Analysis and Preliminary Study}

\section{System Description}

The system is a networked, multi-tier hospital monitoring solution.
At its base sits the \textbf{ESP32 embedded acquisition node}, which occupies the
sensor layer and embedded layer. Its role is to collect physiological and environmental
measurements, apply local signal processing to derive meaningful values (BPM,
SpO\textsubscript{2}), and transmit structured records to a cloud database.

% -------------------------------------------------------
% TODO (Backend / Full-system Team):
% -------------------------------------------------------
\begin{center}
	\textit{[Full system description to be extended with backend and web layers.]}
\end{center}

\section{Use Cases}

The following operational scenarios define the expected system behaviour.

\subsection*{Embedded Node Use Cases}

\begin{description}
	\item[UC-1: Initial Configuration.]
	      A technician connects to the ESP32's Access Point and enter Wi-Fi credentials, Firebase credentials, room ID, and patient ID. The device saves all settings to NVS, serves a confirmation page, and reboots.

	\item[UC-2: Normal Operation.]
	      After boot, the device loads credentials from NVS, connects to the hospital Wi-Fi,
	      authenticates with Firebase, and enters continuous monitoring mode.

	\item[UC-3: Wi-Fi Recovery.]
	      If Wi-Fi is lost during operation (or fails on cold start), the embedded device
	      detects the disconnection and retries \texttt{wifiConnect()} every 10 seconds,
	      indefinitely, without manual intervention.

	\item[UC-4: Remote Reconfiguration.]
	      A technician connects to the AP at any time to update credentials.
	      After saving, the device reboots and applies the new configuration.
\end{description}

% -------------------------------------------------------
% TODO (Backend Team): Add use cases for the backend:
% e.g. UC-6: User login, UC-7: Alert configuration,
% UC-8: Historical data query.
% -------------------------------------------------------
\subsection*{Backend Use Cases}
\begin{center}
	\textit{[To be completed by the backend team.]}
\end{center}

% -------------------------------------------------------
% TODO (Web Dashboard Team): Add use cases for the dashboard:
% e.g. UC-9: Real-time monitoring view, UC-10: Alert history.
% -------------------------------------------------------
\subsection*{Web Dashboard Use Cases}
\begin{center}
	\textit{[To be completed by the web dashboard team.]}
\end{center}

% -------------------------------------------------------
% TODO (AI Team): Add use cases for the AI subsystem:
% e.g. UC-11: Anomaly detection, UC-12: Report generation.
% -------------------------------------------------------
\subsection*{AI Subsystem Use Cases}
\begin{center}
	\textit{[To be completed by the AI team.]}
\end{center}

\section{Inputs and Outputs}

\subsection*{Embedded Node — Inputs}

\begin{itemize}
	\item \textbf{DHT22 (GPIO~4, one-wire):} digital temperature (°C) and relative
	      humidity (\%).
	\item \textbf{MAX30102 (I\textsuperscript{2}C fast mode):} raw infrared
	      and red photodiode FIFO values, used to compute BPM and SpO\textsubscript{2}.
\end{itemize}

\subsection*{Embedded Node — Outputs}

\begin{itemize}
	\item \textbf{Firebase RTDB — \texttt{/rooms/\{roomId\}}:}
	      JSON records containing \texttt{temperature}, \texttt{humidity}, and
	      server-side \texttt{timestamp}, pushed on significant change.
	\item \textbf{Firebase RTDB — \texttt{/patients/\{patientId\}}:}
	      JSON records containing \texttt{heartRate}, \texttt{spO2}, and
	      server-side \texttt{timestamp}, pushed on significant change.
\end{itemize}

% -------------------------------------------------------
% TODO (Backend / Web / AI Teams): Describe inputs and
% outputs for each of your subsystems (API endpoints,
% database queries, model inputs/outputs, etc.).
% -------------------------------------------------------
\subsection*{Backend, Web Dashboard, and AI — Inputs/Outputs}
\begin{center}
	\textit{[To be completed by the backend, web dashboard, and AI teams.]}
\end{center}

\section{Constraints}

\subsection*{Memory Constraints}

The ESP32 has 520\,KiB of internal SRAM.
The Firebase FreeRTOS task is allocated 8\,192 bytes of stack, which must accommodate
the TLS handshake context, async client buffers, and local variables.
The NVS partition stores up to ten string key-value pairs for device configuration.

\subsection*{Latency Constraints}

The Firebase upload interval is \texttt{FIREBASE\_INTERVAL\_MS = 1\,000\,ms},
providing a near-real-time update rate.
The DHT22 enforces a hardware minimum inter-sample interval of 2\,000\,ms.

% -------------------------------------------------------
% TODO (Backend Team): Add latency constraints for the API
% -------------------------------------------------------

\subsection*{Energy Constraints}

The device is assumed to be mains-powered in a fixed hospital room installation.
No active power management (light sleep, deep sleep, or radio duty cycling) is
implemented in the current version.

\subsection*{ESP32 Platform Limitations}

\textbf{Sensor timing sensitivity.}
The MAX30102 heart rate and SpO\textsubscript{2} sensor produces a continuous stream
of readings that must be drained from its internal buffer regularly.
If the program is busy doing something else — such as waiting for a Firebase upload
to complete, which can take several hundred milliseconds or more — the sensor buffer
fills up and new samples are dropped.
This results in the sensor reporting zero or no data, as if no finger were present,
even when one is.

\textbf{Why two cores are necessary.}
The ESP32 contains two independent processing cores that can run tasks simultaneously.
This hardware feature is the direct solution to the timing conflict described above.
By assigning sensor reading exclusively to Core~1 and all Firebase communication
exclusively to Core~0, the two tasks never block each other.
The sensor continues reading at its required rate regardless of how long Firebase
takes to authenticate, upload, or reconnect — and Firebase operations proceed without
being interrupted by sensor activity.
% -------------------------------------------------------
% TODO (AI Team): Add constraints for model size, inference
% latency, and memory footprint if deploying on embedded
% hardware, or cloud compute budget if cloud-based.
% -------------------------------------------------------
\subsection*{AI and Cloud Constraints}
\begin{center}
	\textit{[To be completed by the AI team.]}
\end{center}
